\subsubsection{Large Language Model}
% Nil
The Llama-3.1-8B-Instruct model \cite{llama3.1-8b}—the instruction-tuned
8-billion-parameter variant of Meta’s Llama~3.1 family, released in July~2024—was
selected to benchmark large language model (LLM) inference performance across a
single NVIDIA A100 GPU, a single Intel Gaudi2 HPU, a single NVIDIA H100 GPU, and a
single NVIDIA GH200 Grace Hopper Superchip. It is an open-weight model with a
128K-token context window that demonstrates strong performance across reasoning,
coding, and multilingual tasks. With 8~billion parameters, the model is
sufficiently large to stress accelerator memory bandwidth and compute resources
while remaining small enough to run on a single device on every platform, enabling
direct, controlled hardware comparisons free of tensor-parallel or interconnect
confounds.

The Llama~3.1 8B model has been adopted by Intel as a baseline for evaluating
Gaudi2 platforms \cite{intel2024gaudi2}, and was selected by MLCommons as a large
language model pretraining benchmark in the MLPerf Training suite \cite{mlcommons}.
Because MLPerf Training measures pretraining, it uses the base
(non-instruction-tuned) checkpoint; since this study targets inference serving, we
benchmark the instruction-tuned variant, \texttt{Llama-3.1-8B-Instruct}, which
shares the same architecture and parameter count. This makes it a representative
and well-established workload for accelerator evaluation.

The benchmarking pipeline uses vLLM as the serving engine on the three NVIDIA
platforms (A100, H100, and GH200), where it runs on its CUDA backend over the
PyTorch release matching the installed CUDA toolkit \cite{pytorch}. On Gaudi2, the
vLLM Habana extension could not be brought up reliably on our software stack, so
Gaudi2 is instead benchmarked through Intel's Optimum-Habana inference stack in
lazy mode with HPU graphs—the default, vendor-preferred high-performance inference
path for Gaudi2. The serving engine therefore differs between the NVIDIA platforms
and Gaudi2; rather than eliminating the framework confound we control for it,
running each platform through a well-supported software path and recording the
exact framework and version for every run. Graph capture and compilation are
performed during a warmup phase that is excluded from timing on every platform, so
compilation latency is never counted as inference. On the NVIDIA platforms vLLM
additionally measures the prefill phase (time to first token) separately from
decode (inter-token latency); the Optimum-Habana path reports a single aggregate
throughput rather than a prefill/decode split, so Gaudi2 contributes to the
throughput and energy comparisons but not to the latency-decomposition (TTFT/ITL)
results.

Precision is treated as a first-class, independently controlled variable so that
hardware and numerical-format effects are never conflated. Every platform is
benchmarked in BF16, giving a like-for-like BF16-vs-BF16 comparison that isolates
raw hardware performance at identical precision. The three platforms with native
FP8 tensor hardware—Gaudi2, H100, and GH200 (whose Hopper-architecture GPU shares
the H100's FP8 datapath)—are additionally benchmarked in FP8 (E4M3), giving a
like-for-like FP8-vs-FP8 comparison that isolates the 8-bit datapath. The A100
(Ampere) has no native FP8 tensor cores; rather than introducing an emulated,
non-comparable path, it is reported only in BF16 and therefore appears solely in
the BF16 comparison. Critically, Gaudi2 is reported in both BF16 and FP8, so the
throughput and energy benefit attributable to quantization is shown explicitly
rather than being mistaken for a hardware difference.

The two FP8 paths differ by backend and are documented separately. On Gaudi2, FP8
is produced through the Intel Neural Compressor (INC). A one-time calibration pass
instruments each linear layer and records the maximum absolute activation observed
per tensor. These statistics are then postprocessed into per-tensor FP8 scaling
factors. This step is required because FP8 has a much narrower dynamic range than
BF16, so each tensor must be mapped independently to avoid numerical overflow or
excessive precision loss. During quantized inference, both the weights and the
activations of the linear layers are cast to FP8 (E4M3) and the matrix
multiplications execute in FP8 (W8A8) on the Gaudi2 matrix-multiplication engine
(MME), together with an FP8 KV cache; the language-model head and the non-linear
operations remain in BF16. On the NVIDIA Hopper platforms (H100 and GH200), FP8 is
instead applied online and dynamically by vLLM directly from the BF16 checkpoint,
together with an FP8 KV cache, requiring no separate calibration step or quantized
model artifact. In every configuration the KV-cache precision matches the compute
precision—an FP8 KV cache for the FP8 runs on Gaudi2, H100, and GH200, and a BF16
KV cache for all BF16 runs—so the FP8-vs-FP8 and BF16-vs-BF16 comparisons remain
like-for-like; within FP8 the only backend difference is the quantization method
described above.

All experiments use a fixed workload shape of 512 input tokens and 256 output
tokens, with request concurrency swept over 1, 8, 32, 64, and~128 to trace each
platform from the latency-bound regime to its throughput-saturation knee. For every
configuration, an initial warmup pass is discarded before measurement. On the
NVIDIA (vLLM) platforms each configuration is then repeated three times and we
report the median together with run-to-run variability rather than a single-shot
average, whereas the Optimum-Habana path on Gaudi2 reports throughput aggregated
over its own repeated iterations within a single process. Token throughput and peak
accelerator memory are recorded on every platform; time to first token (TTFT),
inter-token latency or time per output token (ITL/TPOT), and end-to-end latency are
recorded on the NVIDIA platforms, where vLLM exposes the prefill/decode split.
Board power is sampled live during execution via \texttt{nvidia-smi} on the GPUs
and \texttt{hl-smi} on Gaudi2 and integrated into an energy-per-token metric; each
job is allocated so that it sees only its own accelerator, ensuring that no
co-located workload contributes to the sampled board power.
